\section{Motivating Experiment}

Existing state-of-the-art methods that use transformers for decision making train them on sequences consisting of single trajectories~\citep{decision_transformer}. In this section, we show that this setup fails to promote in-context learning of new tasks. Instead, we will demonstrate that \textit{training transformers on sequences of multiple trajectories can enable few-shot learning of unseen tasks}. To illustrate this, we compare the performance of a \textit{single-trajectory} transformer trained on single-trajectory sequences, and a \textit{multi-trajectory} transformer trained on multi-trajectory sequences in a simple setting. 

The single-trajectory transformer's context consists of all past transitions from the current episode $c_t = \{(s_0, a_0, s_1, a_1, \ldots, s_{t-1}, a_{t-1})\}$, and the transformer has to predict the action given the current state and the context $p(a_t | s_t, c_t)$ for each time-step $t$. Hence, given a new task different from the training ones (meaning that we cannot expect much in-weights~\citep{chen2022transdreamer} or zero-shot generalization), the agent can only leverage information from past transitions within the same episode. However, agents typically face different decisions within an episode (\ie{} they rarely see the same state twice), making the problem very challenging.

In contrast, a multi-trajectory transformer's context consists of one or more full trajectories from the same task, as well as all past transitions from the current episode. We will consider the one-shot case in this example where $c_t = \{(s_0^0, a_0^0, \ldots, s_{T}^0, a_{T}^0), (s_0^1, a_0^1, \ldots, s_{t-1}^1, a_{t-1}^1, )\}$. The transformer has to predict the action for the current state $p(a_t^1 | s_t^1, c_t)$ for all $t$.

Hence, when faced with a new task (even if very different from from the training ones), the agent can now leverage information from full trajectories on this task. Note that due to the stochasticity inherent in many environments,  the agent will often be faced with new scenarios that don't appear in its context, so simple memorization is not enough and generalization is required. 
However, since the agent has access to full trajectories, it is less likely that it will encounter states that are too different from the ones in its context so it should perform better.

To verify our hypothesis that \textit{in sequential decision making, multi-trajectory transformers enable better in-context learning than single-trajectory transformers}, we perform the following experiment. We collect expert demonstrations from 100K levels from the \texttt{MiniHack-MultiRoom-N6} task. Using the above protocols, we construct training sequences for the single-trajectory and multi-trajectory transformers and train them using the same hyperparameters\footnote{We set trajectory burstiness $p_b = 1$ for this experiment.}. 
We then evaluate them on the \texttt{MiniHack-Labyrinth} task. To ensure a fair comparison, we condition both transformer variants on a single expert demonstration from the test task. 

\textbf{The multi-trajectory transformer with episodic return of \colorbox{lightblue}{\( 0.780 \pm 0.07 \)}
 outperforms the single-trajectory transformer with \colorbox{lightblue}{\(-0.323 \pm 0.05\)}  by a large margin, indicating the importance of training transformers on sequences of trajectories in order to obtain in-context learning of new sequential decision making tasks (from only a handful of demonstrations and without any weight updates).} For the rest of the results in the paper, we use multi-trajectory transformer as our base model for analysis. 
